%% ma: L10-B12-0068
%% lop: 10
%% chuong: 5
%% bai: 12
%% dang: 10.12.1
%% loai: TLN
%% muc: TH
%% dap_an: "0,6"
%% nguon: "chuyen-de-cac-so-dac-trung-cua-mau-so-lieu-khong-ghep-nhom-toan-10-chuong-trinh-moi.pdf (GV. Phan Nhật Linh) (P01), Bài 01 Phần III Câu 2"
%% kien_thuc: "Lớp 10 Bài 12 (số gần đúng, sai số tuyệt đối, độ chính xác, sai số tương đối, quy tròn số)"
%% trang_thai: chinh-thuc
%% ghi_chu: "Tự duyệt theo yêu cầu giáo viên 22/10/2026."
\begin{ex}
Một phép đo đường kính nhân tế bào cho kết quả là $5\pm0{,}3\ \mu\text{m}$. Đường kính thực của nhân tế bào thuộc một đoạn có độ dài bằng bao nhiêu $\mu$m?
\shortans{0,6}
\loigiai{Đường kính thực thuộc đoạn $[4{,}7;5{,}3]$, có độ dài $5{,}3-4{,}7=0{,}6$.}
\end{ex}
